\documentclass[10pt,aspectratio=169]{beamer}

% --- Codifica e lingua ---
\usepackage[utf8]{inputenc}
\usepackage[italian]{babel}

% --- Pacchetti matematici e grafici ---
\usepackage{amsmath,amssymb,amsfonts,amsthm,mathtools}
\usepackage{bm}
\usepackage{tikz}
\usepackage{tikz-cd}

% --- Stile Beamer essenziale e pulito (ispirato a Oaxaca_2019) ---
\setbeamertemplate{navigation symbols}{}
\useinnertheme{default}
\useoutertheme{default}
\usecolortheme{default}
\usefonttheme{serif}

% Palette cromatica elegante da conferenza/tesi
\definecolor{UniBlue}{RGB}{14,54,115}
\definecolor{UniRed}{RGB}{165,30,30}
\definecolor{UniGreen}{RGB}{20,110,60}
\definecolor{SoftGray}{RGB}{245,247,250}

\setbeamercolor{structure}{fg=UniBlue}
\setbeamercolor{title}{fg=UniBlue}
\setbeamercolor{frametitle}{fg=UniBlue}
\setbeamercolor{block title}{fg=UniBlue,bg=SoftGray}
\setbeamercolor{block body}{bg=SoftGray}
\setbeamercolor{alerted text}{fg=UniRed}

% Impostazione footline sobria con numero di slide
\setbeamertemplate{footline}{
  \leavevmode%
  \hbox{%
  \begin{beamercolorbox}[wd=.75\paperwidth,ht=2.5ex,dp=1.125ex,left,leftskip=1em]{author in head/foot}%
    \usebeamerfont{author in head/foot}\scriptsize\textcolor{gray!80}{Alessandro Ballini \quad|\quad Robustezza di tori invarianti: il teorema twist di Herman}
  \end{beamercolorbox}%
  \begin{beamercolorbox}[wd=.25\paperwidth,ht=2.5ex,dp=1.125ex,right,rightskip=1em]{title in head/foot}%
    \usebeamerfont{title in head/foot}\scriptsize\textcolor{gray!80}{\insertframenumber{} / \inserttotalframenumber}
  \end{beamercolorbox}}%
  \vskip0pt%
}

% --- Macro matematiche coerenti con KAM.tex ---
\newcommand{\T}{\mathbb{T}}
\newcommand{\R}{\mathbb{R}}
\newcommand{\C}{\mathbb{C}}
\newcommand{\Z}{\mathbb{Z}}
\newcommand{\N}{\mathbb{N}}
\newcommand{\eps}{\varepsilon}
\newcommand{\grad}{\nabla}
\newcommand{\id}{\operatorname{id}}
\newcommand{\Id}{\operatorname{Id}}
\newcommand{\abs}[1]{\left|#1\right|}
\newcommand{\norm}[1]{\left\|#1\right\|}
\newcommand{\pa}[1]{\left(#1\right)}
\newcommand{\set}[1]{\left\{#1\right\}}
\newcommand{\blu}[1]{\textcolor{UniBlue}{\textbf{#1}}}
\newcommand{\rosso}[1]{\textcolor{UniRed}{\textbf{#1}}}

% frontespizio
\title{\textbf{Robustezza di tori invarianti lagrangiani quasi-periodici:}\\ \large il teorema \textit{twist} di Herman}
\author{\large Alessandro Ballini}
\institute{\large Università degli Studi di Roma Tor Vergata}
\date{Anno Accademico 2025/2026}

\begin{document}

% slide 1
\begin{frame}[plain]
  \titlepage
\end{frame}

% slide 2
\begin{frame}{Il problema della stabilità nei sistemi dinamici}
  \begin{itemize}
    \item \blu{Origine fisica e celeste:} Lo studio della stabilità a lungo termine ha radici storiche secolari nella meccanica celeste newtoniana (problema degli $n$-corpi e stabilità del Sistema Solare).
    \medskip
    \item \blu{Il teorema di Arnold (1963):}
    \begin{itemize}
      \item Sistema planetario: Sole e $n \ge 2$ pianeti con masse molto più piccole.
      \item Trascurando mutee interazioni, ogni pianeta descrive un'ellisse kepleriana disaccoppiata (sistema \alert{completamente integrabile}).
      \item Per interazioni sufficientemente deboli, nello spazio delle fasi esiste un insieme di \textbf{misura di Lebesgue strettamente positiva} di condizioni iniziali che generano moti \alert{quasi-periodici}.
    \end{itemize}
    \medskip
    \item \blu{La dimostrazione rigorosa di Herman (anni '90):}
    \begin{itemize}
      \item Arnold completò la prova solo per il problema piano dei 3 corpi.
      \item Herman formulò un \textbf{teorema astratto di tori invarianti}, permettendo una chiusura completa e trasparente del problema.
    \end{itemize}
  \end{itemize}
\end{frame}

% slide 3
\begin{frame}{Sistemi completamente integrabili}
  Consideriamo lo spazio delle fasi $T^*\T^n \cong \T^n \times \R^n = \{(\theta, r)\}$ dotato della forma simplettica standard:
  \[
    \omega = \sum_{j=1}^n d\theta_j \wedge dr_j, \qquad X_H = \begin{pmatrix} \partial_r H \\ -\partial_\theta H \end{pmatrix}.
  \]
  \pause
  \begin{block}{Hamiltoniane integrabili in coordinate azione-angolo}
    Un'hamiltoniana $H(\theta, r) = H_0(r)$ dipende unicamente dalle azioni $r$:
    \[
      \begin{cases}
        \dot{\theta} = \partial_r H_0(r) =: \alpha(r) \\
        \dot{r} = 0
      \end{cases}
      \implies 
      \begin{cases}
        \theta(t) = \theta_0 + t\,\alpha(r_0) \pmod{\T^n} \\
        r(t) = r_0
      \end{cases}
    \]
  \end{block}
  \begin{itemize}
    \item Lo spazio delle fasi è \textbf{foliato da tori invarianti lagrangiani} $\T^n \times \{r_0\}$.
    \item Il flusso su ciascun toro è una traslazione lineare a frequenza $\alpha = \alpha(r_0)$.
  \end{itemize}
\end{frame}

% slide 4
\begin{frame}{Sistemi quasi-integrabili e teoria perturbativa classica}
  Consideriamo un sistema perturbato, dipendente formalmente da un parametro $0 \le \eps \ll 1$:
  \[
    H(\theta, r) = H_0(r) + \eps H_1(\theta, r) + \eps^2 H_2(\theta, r) + O(\varepsilon^3)
  \]
  \pause
  \blu{Metodo classico:} cercare un simplettomorfismo generato dal flusso al tempo $1$ di un'hamiltoniana ausiliaria $F$:
  \[
    H \circ \phi_1 = H_0 + \eps \big(H_1 + \{H_0, F\}\big) + O(\eps^2) = H_0 + \eps \big(H_1 - \alpha \cdot \partial_\theta F\big) + O(\eps^2)
  \]
  \pause
  Scomponendo $H_1 = \overline{H}_1(r) + \widetilde{H}_1(\theta, r)$ con $\int_{\T^n} \widetilde{H}_1 d\theta = 0$, si ottiene l'\alert{equazione omologica}:
  \[
    \mathcal{L}_\alpha F = \alpha \cdot \partial_\theta F = \widetilde{H}_1(\theta, r) \implies F(\theta) = \sum_{k \in \Z^n \setminus \{0\}} \frac{c_k}{i(k \cdot \alpha)} e^{i k \cdot \theta}
  \]
 \[
    \inf_{k \neq 0}|k \cdot \alpha| = 0
 \]
\end{frame}

% slide 5
\begin{frame}{Condizione aritmetica diofantea}
  Kolmogorov (1954) intuì che per controllare i piccoli denominatori serviva una stima quantitativa dal basso:
  \begin{block}{Vettori $(\gamma, \tau)$-diofantei}
    Un vettore di frequenza $\alpha \in \R^n$ è detto \blu{$(\gamma, \tau)$-diofanteo} ($D_{\gamma, \tau}$) se:
    \[
      |k \cdot \alpha| \ge \frac{\gamma}{|k|^\tau}, \qquad \forall k \in \Z^n \setminus \{0\} \quad (\gamma > 0, \; \tau > n - 1)
    \]
  \end{block}
  \pause
  \begin{itemize}
    \item \blu{Misura piena:} L'insieme $D_\tau = \bigcup_{\gamma > 0} D_{\gamma, \tau}$ ha \textbf{misura di Lebesgue piena} in $\R^n$.
    \item \blu{Risolubilità dell'equazione omologica:} Per ogni $g$ analitica su $\T^n_{s+\sigma}$ con $\int g = 0$, esiste un'unica soluzione $f$ analitica su $\T^n_s$ soddisfacente:
    \[
      |f|_s \le C \gamma^{-1} \sigma^{-(n+\tau)} |g|_{s+\sigma}
    \]
  \end{itemize}
  \alert{Prezzo da pagare:} inevitabile \textbf{perdita di analiticità} $\sigma > 0$ (arbitraria) ad ogni risoluzione.
\end{frame}

% slide 7
\begin{frame}{Il teorema di Kolmogorov e la superconvergenza}
  \begin{block}{Teorema dei tori invarianti (Kolmogorov, 1954)}
    Fissato $\alpha \in D_{\gamma, \tau}$ e una hamiltoniana imperturbata
    \[
      K^0 = c + \alpha \cdot r + \frac{1}{2}Qr \cdot r + O(r^3), \qquad \det Q \neq 0
    \]
    per ogni perturbazione $H$ sufficientemente vicina a $K^0$, esiste una trasformazione simplettica vicina all'identità che preserva un toro invariante lagrangiano su cui il moto è $\alpha$-quasi-periodico.
  \end{block}
  \pause
  \blu{L'idea rivoluzionaria di Kolmogorov:}
  \begin{itemize}
    \item Non eliminare la perturbazione ordine per ordine ($\eps, \eps^2, \dots$), ma usare uno \textbf{schema iterativo quadratico alla Newton}:
    \[
      \eps_{m+1} \sim \eps_m^2 \qquad (\text{o più precisamente } \eps_m^{\chi}, \; 1 < \chi < 2)
    \]
    \item La \alert{superconvergenza} quadratica schiaccia la perdita geometrica di analiticità $\sigma_m^{-(n+\tau)}$, garantendo la convergenza della successione di trasformazioni simplettiche.
  \end{itemize}
\end{frame}

% slide 8
\begin{frame}{La strategia di Herman: separazione delle difficoltà}
  Negli anni '90, Herman riformula la persistenza dei tori invarianti isolando due passaggi distinti:
  \medskip
  \begin{columns}[t]
    \begin{column}{0.48\textwidth}
      \begin{block}{\blu{Passo 1: Analisi Dura (Hard Analysis)}}
        \begin{itemize}
          \item Fissiamo la frequenza $\alpha \in D_{\gamma, \tau}$.
          \item \textbf{Nessuna ipotesi di non degenerazione} su $K^0$.
          \item Si introduce un \alert{controtermine} $\beta \cdot r$ con $\beta \in \R^n$:
          \[
            H - \beta \cdot r = K \circ G
          \]
          \item Inversione locale dell'operatore non lineare $\phi(K, G, \beta) = K \circ G + \beta \cdot r$ con metodo tipo Nash-Moser.
        \end{itemize}
      \end{block}
    \end{column}
    \pause
    \begin{column}{0.48\textwidth}
      \begin{block}{\blu{Passo 2: Eliminazione del Parametro}}
        \begin{itemize}
          \item Il controtermine $\beta$ rompe la coniugazione dinamica di $H$.
          \item Sotto l'ipotesi $\det Q \neq 0$:
          \[
            \beta = \beta(R), \qquad \det \frac{\partial \beta}{\partial R}(0) \neq 0
          \]
          \item \textbf{Teorema della Funzione Implicita classica} (in dimensione finita) per annullare $\beta$:
          \[
            \beta(R^*) = 0 \implies H_{R^*} = K \circ G
          \]
        \end{itemize}
      \end{block}
    \end{column}
  \end{columns}
\end{frame}

% slide 9
\begin{frame}{Il Teorema \textit{Twist} di Herman}
  Definiamo la classe delle hamiltoniane con toro invariante $\alpha$-quasi-periodico all'origine:
  \[
    \mathcal{K}^\alpha = \left\{ K \in \mathcal{H} : K(\theta, r) = c + \alpha \cdot r + O(r^2) \right\} \implies X_K|_{T^n_0} = \begin{pmatrix} \alpha \\ 0 \end{pmatrix}.
  \]
  \begin{block}{Teorema (Twisted Conjugacy di Herman)}
    Siano $\alpha \in D_{\gamma, \tau}$ e $K^0 \in \mathcal{K}^\alpha$. Per ogni hamiltoniana $H$ sufficientemente vicina a $K^0$ nella topologia analitica, esistono unici:
    \begin{enumerate}
      \item $K \in \mathcal{K}^\alpha$ vicino a $K^0$,
      \item $G$ simplettomorfismo esatto vicino all'identità,
      \item un vettore di controtermine $\beta \in \R^n$ vicino a $0$,
    \end{enumerate}
    tali che:
    \[
      \alert{H - \beta \cdot r = K \circ G}
    \]
    in un intorno del toro $T_0^n = \T^n \times \{0\}$.
  \end{block}
\end{frame}

% slide 10
\begin{frame}{Necessità geometrica del controtermine $\beta$}
  \begin{itemize}
    \item \blu{Invarianza del vettore frequenza:} Se due flussi quasi-periodici su $\T^n$ sono topologicamente coniugati tramite un omeomorfismo omotopo all'identità, le loro frequenze devono coincidere esattamente:
    \[
      \alpha = \beta
    \]
    \item \blu{Non stabilità aperta:} Avere un toro invariante con frequenza fissata $\alpha$ \textbf{non è una proprietà aperta} (cioà stabile per piccole perturbazioni).
  \end{itemize}
\end{frame}

% slide 11
\begin{frame}{Struttura dei simplettomorfismi ammissibili}
  Per preservare la natura lagrangiana del toro, consideriamo lo spazio $\mathcal{G}$ dei germi di simplettomorfismi esatti $G: (\theta, r) \mapsto (\Theta, R)$:
  \[
    G(\theta, r) = \Big(\varphi(\theta), \; {}^t\varphi'(\theta)^{-1}\big(r + \rho(\theta)\big)\Big)
  \]
  \begin{itemize}
    \item $\varphi : \T^n \to \T^n$ è un diffeomorfismo del toro che raddrizza il moto sul toro ($\varphi(0) = 0$).
    \item $\rho(\theta) = dS(\theta)$ è una $1$-forma \textbf{esatta} su $\T^n$ che trasla le azioni per riposizionare il toro nello spazio delle fasi.
  \end{itemize}
  \pause
  \begin{block}{Hamiltoniana 'modificata'}
    Il toro $G^{-1}(T^n_0)$ è invariante per il flusso associato all'hamiltoniana 'modificata' $\widehat{H} = H - \beta \cdot r = K \circ G$.
  \end{block}
\end{frame}

% slide 12
\begin{frame}{Setting funzionale: scale di Banach e norme deformate}
  Per gestire rigorosamente le perdite di analiticità senza restringere brutalmente i domini:
  \begin{itemize}
    \item Dominio complesso: $T^n_s = \{(\theta, r) \in \C^n/2\pi\Z^n \times \C^n : \max(|\mathrm{Im}\,\theta|, |r|) \le s\}$.
    \item Spazio di Banach $(\mathcal{H}_s, |\cdot|_s)$ delle funzioni analitiche con norma del sup su $T^n_s$.
  \end{itemize}
  \pause
  \begin{block}{Norme deformate dalla composizione}
    Dato che $G \in \mathcal{G}$ deforma geometricamente il toro, definiamo la norma su $G^*\mathcal{A}(T^n_s)$:
    \[
      |H|_{G, s} := |H \circ G^{-1}|_s
    \]
    Ciò consente di preservare la struttura di Banach e di misurare l'operatore non lineare:
    \[
      \phi : (K, G, \beta) \in \mathcal{K}_{s+\sigma} \times \mathcal{G}_{s+\sigma}^\sigma \times \R^n \longmapsto H = K \circ G + \beta \cdot r \in \mathcal{H}_s
    \]
  \end{block}
\end{frame}

% slide 13
\begin{frame}{Linearizzazione e stime di derivata prima e seconda}
  Calcolando il differenziale di $\phi(x) = K \circ G + \beta \cdot r$ in $x=(K, G, \beta)$:
  \[
    \phi'(x)\delta x = (K' \circ G)\delta G + \delta K \circ G + \delta\beta \cdot r = \delta H
  \]
  Componendo a destra con $G^{-1}$ e proiettando sui gradi di Taylor in $r$:
  \begin{enumerate}
    \item \blu{Scelta di $\delta\beta$ e $\delta c$:} Assorbono le componenti di media spaziale (termini non oscillanti):
    \[
      \delta\beta = \left(\int_{\T^n} \varphi' \circ \varphi^{-1} d\theta\right)^{-1} \int_{\T^n} \dot{H}_1 \, d\theta \quad \delta c = \int_{\T^n} \dot{H}_0 \, d\theta + \int_{\T^n} \rho \circ \varphi^{-1} \, d\theta \cdot \delta\beta
    \]
    \item \blu{Risoluzione di $\dot{\rho}$ e $\dot{\varphi}$:} Si riducono a \textbf{due equazioni omologiche disaccoppiate}:
    \[
      \mathcal{L}_\alpha f = \widetilde{\dot{H}}_0 \implies \dot{\rho} = -df, \qquad \mathcal{L}_\alpha \dot{\varphi} = 2K_2(\theta)\dot{\rho} + \widetilde{\dot{H}}_1
    \]
  \end{enumerate}
  \begin{alertblock}{Stime fondamentali alla Nash-Moser}
    \[
      |\phi'(x)^{-1}\delta H|_s \le C' \sigma^{-(\tau+n+1)} |\delta H|_{G, s+\sigma}, \qquad |\phi''(x)\delta x^{\otimes 2}|_{G, s} \le C'' \sigma^{-2} |\delta x|_{s+\sigma}^2
    \]
  \end{alertblock}
\end{frame}

% slide 14
\begin{frame}{Schema di Newton a perdita controllata di analiticità}
  Fissiamo le riduzioni di dominio ad ogni passo dello schema di Newton:
  \[
    \sigma_k = \frac{\sigma}{6} 2^{-k}, \qquad s_{k+1} = s_k - 3\sigma_k, \qquad 3\sum_{k=0}^\infty \sigma_k = \sigma > 0
  \]
  \begin{columns}[c]
    \begin{column}{0.55\textwidth}
      \begin{itemize}
        \item Residuo quadratico:
        \[
          |x_{k+1} - x_k|_{s_{k+1}} \le c_k |x_k - x_{k-1}|_{s_k}^2
        \]
        con costante di perdita $c_k = \frac{1}{2} C \sigma_k^{-\tau} \sim 2^{k\tau}$
        \medskip
        \item \blu{La superconvergenza vince la perdita:}
        \[
          |x_{k+1} - x_k|_{s_{k+1}} \le \Big(\eps \prod_{j=0}^k c_j^{2^{-j}}\Big)^{2^k} \le \text{cost} \cdot \eps^{2^k}
        \]
      \end{itemize}
    \end{column}
    \begin{column}{0.42\textwidth}
      \begin{block}{Regolarità di $\psi = \phi^{-1}$}
        \begin{itemize}
          \item La serie $\sum (x_{k+1} - x_k)$ converge nello spazio analitico limite $E_s$.
          \item L'inversa locale $\psi : H \mapsto (K, G, \beta)$ è \textbf{lipschitziana} e di classe $C^1$.
        \end{itemize}
      \end{block}
    \end{column}
  \end{columns}
\end{frame}

% slide 15
\begin{frame}{Ritorno al Teorema di Kolmogorov e Conclusioni}
  \blu{Passo finale: Annullamento del controtermine $\beta$}
  \begin{itemize}
    \item Introduciamo una famiglia di perturbazioni per traslazione delle azioni: $r \mapsto r + R$.
    \item Per l'hamiltoniana traslata imperturbata $K^0(r+R)$, il controtermine vale:
    \[
      \widehat{\beta}^0(R) = 2QR + O(R^2) \implies \left.\frac{\partial \widehat{\beta}^0}{\partial R}\right|_{R=0} = 2Q
    \]
    \item Poiché per ipotesi di Kolmogorov $\det Q \neq 0$ (non degenerazione), $R \mapsto \beta(R)$ è un \textbf{diffeomorfismo locale} per ogni $H$ vicina a $K^0$.
    \item Esiste un unico $R^*$ tale che $\beta(R^*) = 0$:
    \[
      \alert{H(\theta, r+R^*) = K \circ G(\theta, r)}
    \]
    Il toro invariante con moto $\alpha$-quasi-periodico persiste.
  \end{itemize}
\end{frame}

\end{document}